\section{Evaluation}\label{sec:evaluation}


% The source counts were generated by:
% $ cd benchmarks/haskell16/pos/
% $ find . -type f -name '*.hs' -exec sed -i '' s/\{-@/\{-#LH/ {} +
% $ sloccount
% For verified-instances:
% $ find src/Data/VerifiedEq* src/Data/VerifiedOrd* -type f -name '*.hs' -exec sloccount {} \+
% For LVish:
% $ git checkout master ; sloccount src/lvish/Data/LVar/PureSet.hs src/lvish/Data/LVar/SLSet.hs
% $ git checkout verified ; sloccount src/lvish/Data/LVar/PureSet.hs src/lvish/Data/LVar/SLSet.hs
% $ sloccount allpairs.hs ; sloccount allpairs_verified.hs
% $ sloccount IntegerSumReduction2NoVerification.hs ; sloccount IntegerSumReduction2.hs

\begin{figure}[t!]
\begin{minipage}[b]{0.47\textwidth}

%\begin{center}
\setlength{\tabcolsep}{3pt}
\begin{tabular}{lllr}
\toprule
  \multicolumn{3}{l}{\textbf{CATEGORY}}              & \textbf{LOC} \\
\toprule
  \textbf{I.} & \multicolumn{3}{l}{\textbf{Arithmetic}} \\[0.05in]
   & Fibonacci      & \S~\ref{sec:overview}          &  48 \\ % Overview.hs
   & Ackermann      & \citep{ackermann}              & 280 \\ % Ackermann.hs

  \midrule

  \textbf{II.} & \multicolumn{3}{l}{\textbf{Algebraic Data Types}} \\[0.05in]

  & Fold Universal & \citep{agdaequational}          & 105 \\ % FoldrUniversal.hs
%  & Fold Fusion    & \citep{agdaequational}          &     \\

  \midrule

  \textbf{III.} & \textbf{Typeclasses} & Fig~\ref{fig:laws} & \\[0.05in]
  & Monoid         & \tPeano, \tMaybe, \tList        & 189 \\ % Monoid*.hs
  & Functor        & \tMaybe, \tList, \tId, \tReader & 296 \\ % Functor*.hs - FunctorReader.hs
  & Applicative    & \tMaybe, \tList, \tId, \tReader & 578 \\ % Applicative*.hs
  & Monad          & \tMaybe, \tList, \tId, \tReader & 435 \\ % Monad*.hs

  \midrule

  \textbf{IV.} & \multicolumn{3}{l}{\textbf{Functional Correctness}} \\[0.05in]
  & SAT Solver     & \citep{Zombie}                  & 133 \\ % Solver.hs
  & Unification    & \citep{Sjoberg2015}             & 200 \\ % Unification

  \midrule

  \textbf{V.} & \multicolumn{3}{l}{\textbf{Deterministic Parallelism}} \\[0.05in]
  & Conc. Sets     & \S~\ref{sec:set}           & 906 \\ % VerifiedEq/Ord + PureSet.hs + SLSet.hs
  & $n$-body       & \S~\ref{sec:nbody}         & 930 \\ % VerifiedEq/Ord + Inj/Iso + allpairs.hs
  & Par. Reducers  & \S~\ref{sec:reducer}       &  55 \\ % VerifiedSemigroup + Iso + IntegerSumReduction2.hs

  \midrule

  \multicolumn{3}{l}{\textbf{TOTAL}}                 & 4155 \\
\bottomrule
\end{tabular}
%\end{center}
\caption{\textbf{Summary of Case Studies}}
\label{fig:eval-summary}

\end{minipage}
\hspace{0.2in}
\begin{minipage}[b]{0.47\textwidth}

%\begin{figure}[t!]
\begin{center}
%\small
\begin{tabular}{rl}

\toprule

\multicolumn{2}{c}{\textbf{Monoid}} \\
{Left Id.}  & $\emempty\ x\ \emappend\  \equiv x$  \\
{Right Id.} & $x\ \emappend\ \emempty \equiv x$  \\
{Assoc}     & $(x\ \emappend\ y)\ \emappend\ z \equiv x\ \emappend\ (y\ \emappend\ z)$ \\

\midrule

\multicolumn{2}{c}{\textbf{Functor}} \\
{Id.}    & $\efmap\ \eid\ xs \equiv \eid\ xs$ \\
{Distr.} & $\efmap\ (g\ecompose\ h)\ xs \equiv (\efmap\ g\ \ecompose\ \efmap\ h)\ xs$\\

\midrule

\multicolumn{2}{c}{\textbf{Applicative}} \\

{Id.}    & $\epure \eid \eseq\ v \equiv v$ \\
{Comp.}  & $\epure (\ecompose) \eseq u \eseq v \eseq w \equiv u \eseq (v \eseq w)$ \\
{Hom.}   & $\epure\ f\ \eseq\ \epure\ x \equiv \epure\ (f\ x)$\\
{Inter.} & $u\ \eseq\ \epure\ y \equiv \epure\ (\$\ y) \ \eseq \ u$ \\

\midrule
\multicolumn{2}{c}{\textbf{Monad}} \\
{Left Id.}   & $\ereturn\ a \ebind f \equiv f\ a$ \\
{Right Id.}  & $m \ebind \ereturn \equiv m$ \\
{Assoc}         & $(m\ebind f) \ebind g \equiv m\ebind (\lambda x \rightarrow f\ x \ebind g)$\\
\bottomrule
\end{tabular}
\end{center}
\caption{\textbf{Summary of Verified Typeclass Laws}}
\label{fig:laws}
%\end{figure}

\end{minipage}
\end{figure}

We have implemented refinement reflection
in \toolname.
%
In this section, we evaluate our approach
by using \toolname to verify a variety of
deep specifications of Haskell functions
drawn from the literature and categorized
in Figure~\ref{fig:eval-summary},
totalling about 4,000 lines of specifications
and proofs.
%
Verification is fast: from 1 to 7s per
file (module), which is about twice as
much time as GHC takes to compile the
corresponding module.
%
Overall there is about one line of
type specification (theorems)
for three lines of code (implementations and proof).
%
Next, we detail each of the five classes of
specifications, illustrate how they were
verified using refinement reflection, and
discuss the strengths and weaknesses of
our approach.
%
\emph{All} of these proofs require refinement
reflection, \ie are beyond the scope of shallow
refinement typing.

%\RGS{Should we make a mention of VIFP here?}
%\NV{Proof strategies can go away}
\mypara{Proof Strategies.}
%
Our proofs use three building blocks, that are seamlessly
connected via refinement typing:
%
\begin{itemize}
  \item \emphbf{Un/folding}
     definitions of a function @f@ at
     arguments @e1...en@, which due
     to refinement reflection, happens
     whenever the term @f e1 ... en@
     appears in a proof.
     For exposition, we render the function
     whose un/folding is relevant as @#f#@;
%
  \item \emphbf{Lemma Application}
     which is carried out by using
     the ``because'' combinator
     ($\because$) to instantiate
     some fact at some inputs;
%
  \item \emphbf{SMT Reasoning}
     in particular, \emph{arithmetic},
     \emph{ordering} and \emph{congruence closure}
     which kicks in automatically (and predictably!),
     allowing us to simplify proofs by not
     having to specify, \eg which subterms
     to rewrite.
\end{itemize}

\subsection{Arithmetic Properties} \label{subsec:arith} \label{subsec:ackermann}

The first category of theorems pertains to the textbook
Fibonacci and Ackermann functions.
%
In \S~\ref{sec:overview} we proved that @fib@ is increasing.
%
We prove a higher order theorem that lifts increasing functions
to monotonic ones:
%
\begin{mcode}
  fMono :: f:(Nat -> Int) -> fUp:(z:Nat -> {f z <= f (z+1)}) -> x:Nat -> y:{x < y} -> {f x <= f y}
\end{mcode}
%
By instantiating the function argument @f@ of @fMono@ with
@fib@, we proved monotonicity of @fib@.
%
\begin{mcode}
  fibMono :: n:Nat -> m:{n < m} -> {!fib! n <= !fib! m}
  fibMono = fMono fib fibUp
\end{mcode}
%
Using higher order functions like @fMono@,
we mechanized the proofs of the Ackermann function
properties from~\cite{ackermann}.
%
We proved 12 equational and arithmetic properties using various
proof techniques, including
instantiation of higher order theorems (like @fMono@),
proof by construction, and
natural number and generalized induction.

\subsection{Algebraic Data Properties}
\label{subsec:fold}

%The second category of properties pertains to data types.

\mypara{Fold Universality}
%Next, we proved properties of list folding,
as encoded in Adga~\citep{agdaequational}.
%
The following encodes the universal property of @foldr@
%
\begin{code}
  foldr_univ :: f:(a -> b -> b)
             -> h:(L a -> b)
             -> e:b
             -> ys:L a
             -> base:{h [] == e }
             -> step:(x:a -> xs:[a] -> {h (x:xs) == f x (h xs)})
             -> {h ys == foldr f e ys}
\end{code}
%
The \toolname proof term of @foldr_univ@
is similar to the one in Agda.
But, there are two major differences.
%
First, specifications do not support
quantifiers: universal quantification
of @x@ and @xs@ in the step assumption
is encoded as functional arguments.
%
Second, unlike Agda, \toolname does not
support optional arguments, thus to use
@foldr_univ@ the proof arguments @base@
and @step@ that were implicit arguments
in Agda need to be explicitly provided.
%
For example, the following code calls
@foldr_universal@ to prove @foldr_fusion@
by explicitly applying proofs for the
base and the step arguments
%
\begin{code}
  foldr_fusion :: h:(b -> c)
               -> f:(a -> b -> b)
               -> g:(a -> c -> c)
               -> e:b
               -> z:[a]
               -> x:a -> y:b
               -> fuse: {h (f x y) == g x (h y)})
               -> {(h.foldr f e) z == foldr g (h e) z}

  foldr_fusion h f g e ys fuse = foldr_univ g (h . foldr f e) (h e) ys
                                   (fuse_base h f e)
                                   (fuse_step h f e g fuse)
\end{code}
%
Where, for example, @fuse_base@ is a function with type
%
\begin{code}
  fuse_base :: h:(b -> c) -> f:(a -> b -> b) -> e:b -> {(h . foldr f e) [] == h e}
\end{code}

\subsection{Typeclass Laws}

We used \toolname to prove the Monoid, Functor,
Applicative, and Monad Laws, summarized in
Figure~\ref{fig:laws}, for various user-defined
instances summarized in Figure~\ref{fig:eval-summary}.

%% The purpose of these proofs is to investigate the
%% proving abilities of \libname.
%% For this purpose, we defined the appropriate class
%% operators on user defined lists, instead of using
%% Haskell's predefined class instances.
%% %
%% In the near future, we plan to embed these proofs
%% to check the laws on real Haskell instances,
%% but this requires some engineering from the
%% \liquidHaskell team.

\mypara{Monoid Laws}
%
A Monoid is a datatype equipped with an associative
binary operator $\emappend$ (@mappend@) and an \emph{identity}
element $\emempty$.
%
We prove that
%
@Peano@ (with a suitable @add@ and @Z@),
@Maybe@ (with a suitable @mappend@ and @Nothing@), and
@List@ (with append @++@ and @[]@) satisfy the monoid laws.
%
%%For example, we prove that @++@ (\S~\ref{subsec:list})
%%is associative by reifying the textbook proof~\cite{HuttonBook}
%%into a Haskell function, where the induction
%%corresponds to case-splitting and recurring
%%on the first argument:
%%%
%%\begin{mcode}
%%assoc :: xs:[a] -> ys:[a] -> zs:[a] ->
%%       {(xs ++ ys) ++ zs = xs ++ (ys ++ zs)}
%%
%%assoc [] ys zs     = ([] ++ ys) ++ zs
%%                   =. [] ++ (ys ++ zs)
%%                   ** QED
%%assoc (x:xs) ys zs = ((x:  xs)++ ys) ++ zs
%%                   =. (x: (xs ++ ys)) ++ zs
%%                   =.  x:((xs ++ ys) ++ zs)
%%                   =.  x: (xs ++ (ys ++ zs))
%%                       $\because$ assoc xs ys zs
%%                   =. (x:xs)  ++ (ys ++ zs)
%%                   ** QED
%%\end{mcode}
%%

\mypara{Functor Laws}
%
A type is a functor if it has a function
@fmap@ that satisfies the \emph{identity}
and \emph{distribution} (or fusion) laws
in Figure~\ref{fig:laws}.
%
For example, consider the proof of
the @map@ distribution law for the lists,
also known as ``map-fusion'', which is the
basis for important optimizations in
GHC~\cite{ghc-map-fusion}.
%
We reflect the definition of @map@ for lists
%
%%\begin{code}
%%  reflect map :: (a -> b) -> [a] -> [b]
%%  map f []     = []
%%  map f (x:xs) = f x : fmap f xs
%%\end{code}
%
and use it to specify fusion and verify it by an inductive proof:
% by induction (recursion) on the list argument:
%
\begin{mcode}
  map_fusion :: f:(b -> c) -> g:(a -> b) -> xs:[a] -> {map (f . g) xs = (map f . map g) xs}
\end{mcode}

%%\begin{mcode}
%%  map_fusion f g []
%%    =  ((map f) #.# (map g)) []
%%    =. (map f) (#map# g [])
%%    =. #map# f []
%%    =. []
%%    =. #map# (f . g) []
%%    ** QED
%%
%%  map_fusion f g (x:xs)
%%    =   #map# (f . g) (x:xs)
%%    =. (f . g) x : map (f . g) xs
%%    =. (f . g) x : (map f #.# map g) xs
%%       $\because$  map_fusion f g xs
%%    =. (f# . #g) x : map f (map g xs)
%%    =. f   (g  x): map f (map g xs)
%%    =. #map# f (g x: map g xs)
%%    =. map f   (#map# g (x:xs))
%%    =. (map f #.# map g)(x:xs)
%%    ** QED
%%\end{mcode}
%%
% \NV{Say why we need defunctionalization}
% \NV{We use app function like HALO (link to the theory)}

% \mypara{Applicative}
\mypara{Monad Laws}
%
The monad laws, which relate the
properties of the two operators
$\ebind$ and $\ereturn$ (Figure~\ref{fig:laws}),
refer to $\lambda$-functions which
are encoded in our logic as uninterpreted functions.
%%thus their proof exercises
%%our support for defunctionalization
%%and $\eta$- and $\beta$-equivalence.
% and the extensionality axioms to prove.
%
For example, we can encode the associativity property
as a refinement type alias:
%
\begin{code}
  type AssocLaw m f g = { m >>= f >>= g = m >>= (\x -> f x >>= g) }
\end{code}
%
and use it to prove that the list-bind is associative:
%
\begin{mcode}
  assoc :: m:[a] -> f:(a ->[b]) -> g:(b ->[c]) -> AssocLaw m f g
  assoc [] f g     =  [] #>>=# f >>= g
                   =. [] #>>=# g
                   =. []
                   =. [] #>>=# (\x -> f x >>= g) ** QED
  assoc (x:xs) f g =  (x:xs) #>>=# f  >>= g
                   =. (f x ++ xs >>= f) >>= g
                   =. (f x >>= g) ++ (xs >>= f >>= g) $\because$ bind_append (f x) (xs >>= f) g
                   =. (f x >>= g) ++ (xs >>= \y -> f y >>= g) $\because$ assoc xs f g
                   =. (\y -> f y >>= g) x ++ (xs >>= \y -> f y >>= g) $\because$ $\beta$eq f g x
                   =. (x:xs) #>>=# (\y -> f y >>= g) ** QED
\end{mcode}
%
Where the @bind_append@ fusion lemma states that:
%
\begin{code}
bind_append
  :: xs:[a] -> ys:[a] -> f:(a -> [b]) -> {(xs++ys) >>= f = (xs >>= f)++(ys >>= f)}
\end{code}
%
Notice that the last step requires
$\beta$-equivalence on anonymous
functions, which we get by explicitly
inserting the redex in the logic,
via the following lemma with @trivial@ proof
%
\begin{mcode}
    $\beta$eq :: f:_ -> g:_ -> x:_ -> {f x >>= g = (\y -> f y >>= g) x}
    $\beta$eq _ _ _ = trivial
\end{mcode}

\subsection{Functional Correctness} \label{subsec:programs}

Next, we proved correctness of two programs
from the literature: a unification
algorithm and a SAT solver.

\mypara{Unification}
%
We verified the
unification of first order terms, as
presented in~\citep{Sjoberg2015}.
%
First, we define a predicate alias for
when two terms @s@ and @t@ are equal
under a substitution @su@:
%
\begin{code}
  eq_sub su s t = apply su s == apply su t
\end{code}
%
Now, we defined a Haskell function
@unify s t@ that can diverge, or return
@Nothing@, or return a substitution @su@
that makes the terms equal:
%
\begin{code}
  unify :: s:Term -> t:Term -> Maybe {su | eq_sub su s t}
\end{code}
%
For specification and verification,
we only needed to reflect @apply@ and
not @unify@; thus, we only had to verify
that the former terminates, and not the latter.
%
We proved correctness by invoking
separate helper lemmata.
%
For example, to prove the post-condition
when unifying a variable @TVar i@ with
a term @t@ in which @i@ \emph{does not}
appear, we apply a lemma @not_in@:
%
\begin{mcode}
  unify (TVar i) t2 | not (i Set_mem freeVars t2) = Just (const [(i, t2)] $\because$ not_in i t2)
\end{mcode}
%%
\ie if @i@ is not free in @t@,
the singleton substitution yields @t@:
%
\begin{code}
  not_in :: i:Int -> t:{Term | not (i Set_mem freeVars t)} -> {eq_sub [(i,t)] (TVar i) t}
\end{code}
%
This example highlights the benefits of
partial verification on a legacy programming language:
potential diverging code (\eg the function @unify@) coexists and invokes proof terms.

\mypara{SAT Solver}
%
As another example, we implemented and verified the simple
SAT solver used to illustrate and evaluate
the features of the dependently typed language
Zombie~\citep{Zombie}.
%
The solver takes as input a formula @f@
and returns an assignment that
\emph{satisfies} @f@ if one exists,
as specified below.
%
\begin{code}
  solve :: f:Formula -> Maybe {a:Assignment | sat a f}
\end{code}
%
The function @sat a f@ returns @True@ iff the assignment
@a@ satisfies the formula @f@.
%
Verifying @solve@ follows directly by
reflecting @sat@ into the refinement logic.
